% TOP_PROBLEM: {"id": 1252, "title": "The Number of Idoneal Numbers", "category_id": 1, "category": {"id": 1, "name": "number_theory", "display_name": "Number Theory", "description": "Properties of integers, prime numbers, Diophantine equations.", "slug": "number-theory", "order_index": 1, "created_at": "2026-07-31T15:26:25.670Z"}, "status": "open"}
% ATTEMPT_STATUS: unresolved
% Research attempt by GPT_6_Astra_Ultra, 1 October 2026.
% Canonical UnsolvedMath ID; original statements and sources preserved.
% FIRST_POSTED: 2026-10-01
% Imported from: unsolvedmath_additional_500.tex; UnsolvedMath ID 1252
% !TeX program = lualatex
% Compile twice: lualatex 1252.tex
% Self-contained: no external bibliography, images, or \input files.
% Numeric section labels and visible numbers are original UnsolvedMath IDs.
\documentclass[11pt,a4paper]{article}
\usepackage[margin=24mm,headheight=14pt]{geometry}
\usepackage{fontspec}
\setmainfont{Latin Modern Roman}
\setsansfont{Latin Modern Sans}
\setmonofont{Latin Modern Mono}
\usepackage{amsmath,amssymb,mathtools}
\usepackage{unicode-math}
\setmathfont{Latin Modern Math}
\usepackage{microtype}
\usepackage{enumitem,xurl,needspace}
\usepackage[dvipsnames]{xcolor}
\usepackage{hyperref}
\hypersetup{unicode=true,colorlinks=true,linkcolor=MidnightBlue,urlcolor=MidnightBlue,
 pdftitle={UnsolvedMath Problem 1252},pdfauthor={GPT 6 Astra Ultra},bookmarksnumbered=true}
\usepackage{bookmark,tocloft,titlesec,fancyhdr}
\setlength{\cftsecnumwidth}{4.2em}
\cftsetpnumwidth{2.5em}
\cftsetrmarg{4em}
\titleformat{\section}{\Large\bfseries\raggedright}{\thesection}{0.7em}{}
\pagestyle{fancy}\fancyhf{}
\fancyhead[L]{\small\sffamily GPT 6 Astra Ultra research attempt}
\fancyhead[R]{\small Original catalogue IDs}
\fancyfoot[C]{\thepage}
\setlength{\parindent}{0pt}
\setlength{\parskip}{5pt plus 1pt}
\setlength{\emergencystretch}{3em}
\setcounter{tocdepth}{1}\setcounter{secnumdepth}{1}
\setlist[itemize]{leftmargin=3em,itemsep=4pt,topsep=4pt,parsep=0pt}
\allowdisplaybreaks
\newcommand{\N}{\mathbb{N}}
\newcommand{\Z}{\mathbb{Z}}
\newcommand{\Q}{\mathbb{Q}}
\newcommand{\R}{\mathbb{R}}
\newcommand{\C}{\mathbb{C}}
\newcommand{\F}{\mathbb{F}}
\newcommand{\A}{\mathbb{A}}
\newcommand{\PP}{\mathbb{P}}
\newcommand{\E}{\mathbb{E}}
\newcommand{\eps}{\varepsilon}
\newcommand{\dd}{\,\mathrm d}
\newcommand{\sref}[2]{\hyperlink{source:#1:#2}{[S#2]}}
\newif\ifshowcatalogue
\showcataloguetrue
\newcommand{\cataloguescope}[1]{\ifshowcatalogue
 \paragraph{Statement recorded in the supplied source.}
 #1\fi}
\begin{document}
% UnsolvedMath ID 1252; source catalogue code NT-045

\setcounter{section}{1251}

\section{Completeness of the List of Idoneal Numbers}\label{1252}

{\small\sffamily UnsolvedMath ID 1252 · Number Theory}

\subsection{Definitions and mathematical statement}

\paragraph{Shared notation.}Unless a section says otherwise, $\N=\{1,2,\ldots\}$,
$\Z$ is the ring of integers, $\Q,\R,\C$ are the rational, real, and complex
fields, $[N]=\{1,\ldots,\lfloor N\rfloor\}$, and $\log$ is the natural
logarithm. For real functions of a parameter tending to infinity,
$f\ll g$ and $f=O(g)$ mean $|f|\leq Cg$ eventually for some fixed $C$;
$f\gg g$ reverses this comparison; $f=o(g)$ means $f/g\to0$;
$f\sim g$ means $f/g\to1$; and $f\asymp g$ means both order bounds.
A subscript on $O$ or $\ll$ specifies allowed constant dependence.
The expression $n^{o(1)}$ in an upper bound means growth smaller than
$n^\epsilon$ for each fixed $\epsilon>0$. Local definitions govern any
exception, finite-field convention, or use of the same symbol for another
object. An asymptotic claim is not automatically an effective algorithm.



An idoneal positive integer $n$ can be characterized by the proper ideal class group of the imaginary quadratic order of discriminant $-4n$: every ideal class has order dividing two. Equivalently one uses proper classes of primitive positive definite binary quadratic forms of that discriminant. The target is the completeness of the classical list, not merely a search up to a fixed bound. \sref{1252}{1} \sref{1252}{2}

\paragraph{Statement recorded in the supplied source.}
Are there exactly 65 idoneal numbers, or could there be 66 or 67? \sref{1252}{1}

\paragraph{Residual task reported by the archive.}

The embedded review (2026-08-17) records: Exclude or exhibit the one or two possible exceptional idoneal numbers without assuming GRH. \sref{1252}{1}

\subsection{Short English statement}

Is the classical collection of sixty-five idoneal integers complete, or do further examples exist?

\subsection{Research attempt}

\paragraph{Scope and approach.}
The class group in the definition belongs to the order of discriminant
$-4n$, which need not be the maximal order of its quadratic field.
Replacing it by the maximal-order class group could change the answer.
Kani's primary survey [E1] treats this distinction and corrections to
historical formulations. This attempt uses the stated order throughout
and constructs an exact finite test by reduced forms. It does not assume
the generalized Riemann hypothesis or an effective global cutoff.

\paragraph{A complete finite test for a fixed discriminant.}
Write a primitive positive definite form as $(a,b,c)$ with
$b^2-4ac=-4n$. Use the usual reduced representative convention
\[
 |b|\le a\le c,\quad \gcd(a,b,c)=1,
 \quad b\ge0\text{ if }|b|=a\text{ or }a=c.
\]
The reduction theorem for positive definite forms gives one such
representative per proper class under this boundary convention. This
standard theorem is an input from quadratic-form theory; it is not
reproved here. The bound needed for enumeration follows immediately:
\[
 4n=4ac-b^2\ge4a^2-a^2=3a^2,
 \qquad a\le\sqrt{4n/3}.
\]
Thus enumerate those finitely many positive $a$, the integers
$-a\le b\le a$, and retain $c=(b^2+4n)/(4a)$ when integral and
satisfying the displayed conditions. This proves coverage at fixed $n$.

\paragraph{Detecting exponent two without a composition table.}
The inverse of the proper class represented by $(a,b,c)$ is represented
by $(a,-b,c)$. In the interior $0<|b|<a<c$, these are distinct reduced
representatives, so the class cannot equal its inverse. On the boundary,
the cases $b=0$, $|b|=a$, or $a=c$ are exactly the ambiguous reduced
forms, whose proper classes equal their inverses. For example the
integral substitutions changing a variable by the other, or interchanging
the two variables with a sign, identify the boundary inverse with its
chosen representative. These substitutions have determinant one.
Consequently the required class-group exponent condition is equivalent
to every enumerated form satisfying at least one of
\[
                       b=0,\qquad |b|=a,\qquad a=c.
\]
The test uses proper classes and their inverses, not improper equivalence,
which would identify inverse pairs and lose the obstruction.

\paragraph{Small complete examples.}
For $n=5$ the reduced forms are $(1,0,5)$ and $(2,2,3)$; both are
ambiguous. For $n=6$ they are $(1,0,6)$ and $(2,0,3)$, again both
ambiguous. For $n=11$ there are three forms
\[
                   (1,0,11),\quad (3,-2,4),\quad(3,2,4).
\]
The latter two are distinct inverse classes in the strict interior,
so eleven is not idoneal under the catalogue criterion. These examples
illustrate an actual obstruction, rather than merely checking that a
class number is a power of two. A group of order a power of two may
contain elements of order four or higher, so that numerical class-number
condition alone is insufficient.

\paragraph{What finite reduction does and does not bound.}
The inequality $a\le\sqrt{4n/3}$ bounds representatives after $n$
has been fixed; it puts no bound on $n$. To prove completeness of the
classical list by this route one needs to show that every larger
discriminant has a primitive interior reduced form with nonzero $b$,
or obtain a separate effective discriminant bound. An algorithm which
always halts for each individual $n$ can still run forever over all
$n$ when no further idoneal value exists. Thus a decision procedure
at fixed input is different from a proof that no exceptional input
remains.

One possible direct obstruction is to exhibit $a,b$ with
$0<|b|<a<c=(b^2+4n)/(4a)$ and $\gcd(a,b,c)=1$. The divisibility
$4a\mid b^2+4n$ is a local condition and primitivity is another.
Producing such witnesses uniformly in all sufficiently large $n$ is
precisely the unproved step of this elementary approach. Selecting $a$
near $\sqrt n$ by an approximation does not ensure this congruence,
and discarding nonmaximal orders does not preserve the target.

\paragraph{Executed finite audit.}
The integer script enumerates all primitive reduced forms for every
$1\le n\le5000$ and applies the inverse-class criterion. Exactly
$65$ values pass; the largest is $1848$. The full list and sample form
sets, including all sixteen forms for $n=1848$, are in
\texttt{results\_second.json}. The code is
\texttt{checks\_second.py}, both under
\path{_Local/Erdos_problems/UnsolvedMath/attempt_audit_2026_10_01/number_theory/}.
The result is a complete classification only within the stated bound.
No interval beyond $5000$ has been excluded by this calculation.

\paragraph{Remaining gap and outcome.}
The attempt supplies exact inverse-class witnesses and a reproducible
finite test respecting the order discriminant. It supplies no effective
global bound and no uniform construction of a nonambiguous class at all
larger discriminants. \textbf{Outcome: UNRESOLVED.} The finite count of
$65$ is not promoted to an unconditional proof of completeness.

\subsection{Sources}

{\small

\begin{itemize}

\item[\hypertarget{source:1252:1}{S1}] ULAM.AI, \emph{UnsolvedMath}, supplied \texttt{problems.json}, record 1252 (NT-045), statement and background. The embedded literature-review date is 2026-08-17; that is the archive’s review date, not a fresh certification. Dataset landing page: \url{https://huggingface.co/datasets/ulamai/UnsolvedMath}.

\item[\hypertarget{source:1252:2}{S2}] Peter J. Weinberger, Exponents of the class groups of complex quadratic fields, Acta Arithmetica 22 (1973), 117-124. \url{https://eudml.org/doc/205250}. Bibliographic information imported from the supplied record.

\item[\hypertarget{source:1252:3}{S3}] Ernst Kani, Idoneal Numbers and Some Generalizations, Annales des sciences mathematiques du Quebec 35 (2011), 197-227. \url{https://mast.queensu.ca/~kani/papers/2011K1.pdf}. Bibliographic information imported from the supplied record.


\item[E1] Ernst Kani, \emph{Idoneal Numbers and some Generalizations},
author manuscript. \url{https://mast.queensu.ca/~kani/papers/idoneal-f.pdf}.
Author-manuscript formulation and search-visible text inspected
1 October 2026; direct PDF retrieval was unavailable. The standard
reduction and inverse-class facts are explicitly identified as inputs.

\end{itemize}

}

\label{end:1252}
\end{document}
